% CV de David Vicente, mis en forme comme un article scientifique :
% classe article standard, Computer Modern (Latin Modern), titre centré,
% résumé et mots-clés, sections numérotées, réalisations en bibliographie
% citées dans le texte.
%
% Compilation : latexmk -pdf cv.tex (deux passes pour les citations),
% ou npm run cv depuis la racine du site, qui copie le PDF dans public/.
%
% Le contenu suit src/data/ du site. À mettre à jour en même temps.

\documentclass[10pt,a4paper]{article}

\usepackage[utf8]{inputenc}
\usepackage[T1]{fontenc}
\usepackage{lmodern}
\usepackage[french]{babel}
\usepackage{microtype}
\usepackage[a4paper,top=13mm,bottom=13mm,left=19mm,right=19mm]{geometry}
\usepackage{graphicx}
\usepackage{enumitem}
\usepackage{cite}
\usepackage[hidelinks]{hyperref}

% Texte extrait du PDF (copie, logiciels de recrutement) : accents conservés.
\input{glyphtounicode}
\pdfgentounicode=1

\hypersetup{
  pdftitle={David Vicente – CV – Ingénieur logiciel, développement web et mobile},
  pdfauthor={David Vicente},
  pdfsubject={Curriculum vitae},
  pdfkeywords={ingénieur logiciel, développement web, développement mobile,
    TypeScript, React, Next.js, React Native, Node.js, Toulouse},
  pdflang={fr-FR},
}

\addto\captionsfrench{\renewcommand{\refname}{Réalisations}}

% Titres de section de la classe article, avec moins d'espace autour, comme
% dans les gabarits de conférence.
\makeatletter
\renewcommand\section{\@startsection{section}{1}{\z@}%
  {-2.6ex \@plus -.6ex \@minus -.2ex}{1.4ex \@plus .2ex}%
  {\normalfont\Large\bfseries}}
\makeatother
\pagestyle{empty}
\setlist{itemsep=0pt, topsep=2pt}

% Intitulé d'un poste ou d'un diplôme, la période en italique à droite.
\newcommand{\entry}[2]{\par\noindent #1\hfill\textit{#2}\par}

\begin{document}

% ----------------------------------------------------------------- En-tête --
% Photo à gauche ; à sa droite le nom, le titre et la ville, alignés à gauche ;
% les coordonnées en colonne, alignées sur la marge droite.
\noindent
\begin{minipage}[c]{23mm}
  \includegraphics[width=\linewidth]{portrait.jpg}
\end{minipage}%
\hspace{6mm}%
\begin{minipage}[c]{\dimexpr\textwidth-23mm-6mm-58mm\relax}
  \raggedright
  {\LARGE David Vicente\par}
  \vspace{5pt}
  {\large Ingénieur logiciel\par}
  \vspace{2pt}
  Développement web et mobile\par
  \vspace{2pt}
  Toulouse, France
\end{minipage}%
\begin{minipage}[c]{58mm}
  \raggedleft
  % +33\,6\,00\,00\,00\,00\\[2pt]  % téléphone, à décommenter
  \href{mailto:david.vct@proton.me}{\texttt{david.vct@proton.me}}\\[2pt]
  \href{https://david-vct.github.io}{\texttt{david-vct.github.io}}\\[2pt]
  \href{https://github.com/david-vct}{\texttt{github.com/david-vct}}\\[2pt]
  \href{https://www.linkedin.com/in/david-vct/}{\texttt{linkedin.com/in/david-vct}}
\end{minipage}

\vspace{4pt}

% ----------------------------------------------------------------- Résumé ---
\begin{abstract}
\noindent
Ingénieur logiciel formé aux universités de Bordeaux et de Toulouse, je conçois
et développe des applications web et mobiles depuis 2021, principalement dans
l’écosystème TypeScript : React, Next.js et React Native côté client, Node.js
côté serveur. Indépendant depuis 2023, j’accompagne PME et startups, aussi bien
pour concevoir un produit de bout en bout~\cite{verse,padelgo,trik3r} que pour
renforcer une équipe existante, et je publie des projets
open source~\cite{experimental,ddgct,ydotool,vocabulary}. J’ai notamment
travaillé pour Capgemini, EDF, le Groupe Adonis, SEDRAP et PadelGo.

\medskip\noindent
\textbf{Mots-clés :} développement web et mobile, TypeScript, React, Next.js,
React Native, Node.js.
\end{abstract}

% ------------------------------------------------------------- Expérience ---
\section{Expérience professionnelle}

\entry{\textbf{Développeur web et mobile}, indépendant, Toulouse}
      {sept.~2023 -- aujourd’hui}
\begin{itemize}
  \item Conception et développement de produits web et mobiles pour des PME et
        des startups, de la maquette à la publication sur les stores :
        Verse~\cite{verse}, PadelGo et PadelGo Manager~\cite{padelgo},
        Trik3r Collecte~\cite{trik3r}.
  \item Renfort technique au sein d’équipes de développement existantes.
  \item Environnement : TypeScript, React, Next.js, React Native, Expo, Node.js.
\end{itemize}

\smallskip
\entry{\textbf{Développeur full stack}, Capgemini, alternance, Toulouse}
      {sept.~2022 -- août~2023}
\begin{itemize}
  \item Conception et implémentation de solutions complètes pour un grand compte.
  \item Coordination et organisation agile selon le cadre SAFe.
  \item Environnement : Java, JavaScript.
\end{itemize}

\smallskip
\entry{\textbf{Développeur full stack}, Groupe Adonis, alternance}
      {sept.~2021 -- sept.~2022}
\begin{itemize}
  \item Conception, design et implémentation d’un tableau de bord.
  \item Développement d’une API REST.
  \item Environnement : HTML, API REST.
\end{itemize}

% ------------------------------------------------------------ Compétences ---
\section{Compétences}

\begin{description}[itemsep=0pt, topsep=2pt, font=\normalfont\bfseries]
  \item[Langages :] TypeScript, JavaScript, Java, SQL, Bash, HTML, CSS.
  \item[Web et mobile :] React, Next.js, React Native, Expo, Tailwind CSS.
  \item[Serveur et données :] Node.js, NestJS, API REST, PostgreSQL, MongoDB,
        Firebase.
  \item[Outils :] Git, Docker, Linux, Grafana, Google Cloud Platform.
  \item[Méthodes :] conception logicielle, gestion de projet, agilité
        (Scrum, SAFe).
  % \item[Langues :] français (langue maternelle), anglais (niveau à préciser).
\end{description}

% -------------------------------------------------------------- Formation ---
\section{Formation}

\entry{\textbf{Master Informatique} (bac\,+\,5), Université de Toulouse, en alternance}
      {2021 -- 2023}
\entry{\textbf{Licence Informatique} (bac\,+\,3), Université de Bordeaux}
      {2018 -- 2021}
\entry{\textbf{Baccalauréat scientifique}, spécialité informatique,
       Lycée Grand Air, Arcachon}
      {2015 -- 2018}

% ----------------------------------------------------------- Réalisations ---
% Produits conçus pour des clients, puis projets personnels open source.
\begin{small}
\begin{thebibliography}{9}
\setlength{\itemsep}{1pt}

\bibitem{verse}
\textit{Verse}. Application mobile de suivi de films, séries, livres et jeux
vidéo, réunis dans une bibliothèque d’univers fictifs : évaluations,
collections et fonctionnalités sociales. React Native, Node.js. Depuis 2026.

\bibitem{padelgo}
\textit{PadelGo} et \textit{PadelGo Manager}. Conception et développement
intégral d’une application iOS et Android de réservation de terrains et de
recherche de partenaires de padel, et de la plateforme web de gestion des clubs.
React Native, Expo, Next.js, TypeScript. 2024 -- 2026.
Disponible sur
l’\href{https://apps.apple.com/ma/app/padelgo/id6743053827?l=fr-FR}{App Store}
et \href{https://play.google.com/store/apps/details?id=com.rmarsile.padelgo}{Google Play}.

\bibitem{trik3r}
\textit{Trik3r Collecte}. Plateforme web de gestion des collectes de déchets :
commandes, planification des interventions et suivi des prestations.
Node.js, TypeScript. 2024.

\bibitem{experimental}
\textit{Experimental Coding}. Galerie de code créatif : quines, fractales,
shaders WebGL et automates cellulaires, chacun en un seul fichier sans
dépendance. JavaScript, WebGL. 2026.
\href{https://david-vct.github.io/experimental-coding/}{\nolinkurl{david-vct.github.io/experimental-coding}}.

\bibitem{ddgct}
\textit{DuckDuckGo Custom Tabs}. Extension de navigateur, publiée sur
\href{https://addons.mozilla.org/en-US/firefox/addon/duckduckgo-custom-tabs/}{Firefox Add-ons},
qui redirige les onglets de DuckDuckGo vers les services de son choix.
TypeScript, WebExtensions. 2026.
\href{https://github.com/david-vct/duckduckgo-custom-tabs}{\nolinkurl{github.com/david-vct/duckduckgo-custom-tabs}}.

\bibitem{ydotool}
\textit{Ydotool Rebind}. Outil en ligne de commande pour Linux : surcouche Bash
de ydotool pour les claviers non QWERTY. 2026.
\href{https://github.com/david-vct/ydotool-rebind}{\nolinkurl{github.com/david-vct/ydotool-rebind}}.

\bibitem{vocabulary}
\textit{Vocabulary Tools for Google Docs}. Module complémentaire donnant accès
aux synonymes, antonymes et conjugaisons depuis l’éditeur. Google Apps Script,
Google Cloud Platform. 2024.
\href{https://github.com/david-vct/gdocs-vocabulary-tools}{\nolinkurl{github.com/david-vct/gdocs-vocabulary-tools}}.

\end{thebibliography}
\end{small}

\end{document}
